% Limite do modelo: 180 palavras cada; sem citações; 3 a 5 palavras-chave em ordem alfabética.

\begin{Abstract}
\textbf{Background:} turning a soil analysis into liming and fertilization recommendations
requires applying dozens of tables from the regional manual, and no tool available to technicians
in southern Brazil combines that recommendation with an assessment of the area's suitability. This
work developed SIRAS, a rule-based expert system that recommends liming and fertilization for 61
crops following the Rio Grande do Sul and Santa Catarina manual and estimates edaphic suitability
before and after correction, citing the rule and table behind every value.
\textbf{Results:} all 25 validation cases matched values calculated by hand from the manual (95\%
confidence interval: 86.7--100\%); the suitability module conformed to its frozen specification in
all 92 cases; and the system reproduced the 32 classifications published in the manual's annex.
Validation also revealed two defects missed by automated tests.
In a usability test, all eight participants completed a recommendation scenario unaided, and no
heuristic was rated above a cosmetic problem.
\textbf{Conclusions:} the system faithfully implements the manual's criteria and integrates
recommendation and suitability in a traceable report; the suitability criteria still require
independent agronomic validation.
\end{Abstract}

\begin{keywords}
Expert systems; liming and fertilization; rule-based systems; soil analysis; soil suitability
\end{keywords}

\begin{resumo}
\textbf{Contexto:} transformar a análise de solo em recomendação de calagem e adubação exige
aplicar dezenas de tabelas do manual regional, e nenhuma ferramenta acessível ao técnico gaúcho
integra essa recomendação à avaliação da aptidão da área. Este trabalho desenvolveu o SIRAS, um
sistema especialista baseado em regras que recomenda calagem e adubação para 61 culturas segundo o
Manual do RS e de SC e estima a aptidão edáfica antes e depois da correção, indicando a regra e a
tabela de cada valor.
\textbf{Resultados:} os 25 casos de validação concordaram com valores calculados à mão (IC 95\%:
86,7--100\%); a aptidão foi conforme à especificação nos 92 casos; e o sistema reproduziu as 32
classificações publicadas no anexo do Manual. A validação revelou dois defeitos que os testes automatizados não detectaram.
No teste de usabilidade, os oito participantes concluíram sozinhos um cenário de recomendação, e
nenhuma heurística recebeu nota acima de problema estético.
\textbf{Conclusões:} o sistema implementa fielmente os critérios do Manual e integra recomendação
e aptidão num laudo rastreável; os critérios de aptidão ainda requerem validação agronômica
independente.
\end{resumo}

\begin{palavras_chave}
Análise de solo; aptidão edáfica; calagem e adubação; sistemas baseados em regras; sistemas especialistas
\end{palavras_chave}
